\documentclass[11pt]{article}
\usepackage[a4paper,margin=29mm]{geometry}
\usepackage{amsmath,amssymb,amsthm,mathtools,booktabs,microtype}
\usepackage[T1]{fontenc}
\usepackage{lmodern}
\usepackage[hidelinks]{hyperref}
\hypersetup{pdftitle={The radical-square invariant for geometric-progression-free polynomial sets},pdfauthor={Research note for ProveIt}}
\newtheorem{theorem}{Theorem}[section]
\newtheorem{lemma}[theorem]{Lemma}
\newtheorem{corollary}[theorem]{Corollary}
\newtheorem{proposition}[theorem]{Proposition}
\theoremstyle{remark}\newtheorem{remark}[theorem]{Remark}
\newcommand{\F}{\mathbb F}
\newcommand{\m}{\mathfrak m}
\newcommand{\E}{\mathbb E}
\newcommand{\Prob}{\mathbb P}
\newcommand{\rad}{\operatorname{rad}}
\newcommand{\dbar}{\overline{d}}
\newcommand{\degbar}{\deg\overline f}
\newcommand{\one}{\mathbf 1}
\title{The radical-square invariant for geometric-progression-free polynomial sets\\\large A product-density classification beyond square-zero local rings}
\author{Research note for ProveIt}
\date{27 November 2026}
\begin{document}
\maketitle
\begin{abstract}
Let $A$ be a finite commutative local ring with maximal ideal $\m$, residue field of order $q$ and order $Q$. Assume $q$ is even, $2\m=0$, and $\m$ annihilates the squares of its elements. Put $S=\{a^2:a\in\m\}$ and $w=|S|/Q^2$. We characterize exactly the polynomials with nonzero reduction that can begin a square-ratio pair within a degree ball. For $\m^3=0$ these are exactly the possible starts of three-term geometric progressions in that ball. A disjoint packing of triples and a deletion construction give $(1-q^{-1})^2w\le 1-\Delta_L(A)\le w$ for every fixed bound $L$ on the sum of the two exponent gaps. For arbitrary growing products of such rings, the limiting maximum upper density is positive if and only if the sum of the $w$'s is finite. Consequently two sequences of rings with identical residue fields, cardinalities, and all radical-layer dimensions can have respectively positive and zero limiting extremal density. This resolves the equicharacteristic-two cube-zero case, including its growing-product criterion, of an open question in the merged ProveIt article. The results are unconditional; no extension of Dirichlet GRH to nonabelian $L$-functions is asserted.
\end{abstract}

\section{Source, target, and conventions}
Asada, Fourakis, Manski, McNew, Miller and Moreland \cite{asada} ask in their Question 6.5 what becomes of geometric-progression avoidance when the coefficient field is replaced by a finite ring. The merged ProveIt article \cite{proveit} develops this question for local rings with square-zero maximal ideal, and asks at \texttt{gpp:q:cube} which invariant replaces the square-zero density-loss parameter when $\m^3=0\ne\m^2$. It also asks for a criterion for arbitrary growing products. We answer that question for all characteristic-two local rings with cube-zero radical, and for a larger class (including some mixed-characteristic rings). The unrestricted odd-characteristic and mixed-characteristic cases are not settled here.

\paragraph{Repository analysis.} The source tree describes an article merging 209 research reports, not an independent verification of all the reports. The vendored \texttt{Statements.lean} and \texttt{Main.lean} explicitly quantify over Dirichlet characters over $\mathbb Q$ and export RH and the Dirichlet zero-free half-plane and critical-line statements. Section \texttt{src.tex} explains the source mechanism: a one-sided smoothed prime bound is obtained by comparing a short-form Gram determinant and an integrality lower bound, with a spanning estimate controlling missing dimensions; a twist separates zero multiplicities and a Landau argument excludes poles. The density chapter transfers the finite selection/packing and elimination aspects rather than the analytic zero criterion. Its theorem \texttt{gpp:thm:local} uses degree elimination in one factor, disjoint triple packings, and the density Hales--Jewett theorem to classify products of square-zero local rings. In even residue characteristic its parameter is $Q^{-2}$; the new parameter below equals this when the squaring map vanishes, but differs when it has a large image. The present argument replaces the square-zero elimination by an exact description of the top two coefficients, then builds a recoverable packing. This is a direct extension of the density chapter, not an attempted derivation of nonabelian GRH from Dirichlet GRH.

\paragraph{Definitions.} Throughout, rings are finite, commutative and unital, with nonzero residue fields. For a finite ring $R$ define
\[
 B_n(R)=\{f\in R[x]:\deg f\le n\},\quad \deg 0=-\infty,
 \qquad |B_n(R)|=|R|^{n+1}.
\]
For $L\ge2$ fixed, a subset is admissible if it contains no three \emph{distinct} polynomials
\begin{equation}\label{eq:pattern}
 f,\quad r^af,\quad r^{a+b}f,\qquad
 r\in R[x]\text{ a nonunit},\qquad a,b\ge1,\quad a+b\le L.
\end{equation}
Let $M_{L,n}(R)$ be the maximum admissible fraction of $B_n(R)$ and let
\[
 \Delta_L(R)=\sup_{E\text{ admissible}}\limsup_{n\to\infty}
 \frac{|E\cap B_n(R)|}{|B_n(R)|}.
\]
Write $\Delta=\Delta_2$. The nonunit-ratio convention is the one in the cited question and the repository. The lower bounds we give also avoid repeated-term patterns with nonunit ratio, because all residue components are nonzero and a nonunit ratio then either forces a zero residue or strictly increases a residue degree. The distinct-term convention is used in all upper bounds.

\section{The exact local obstruction}
Call a finite local ring $A$ \emph{eligible} if, for its maximal ideal $\m$ and residue field $k=A/\m$ of order $q$,
\begin{equation}\label{eq:hyp}
 q\text{ is even},\qquad 2\m=0,\qquad
 \m a^2=0\quad(a\in\m).
\end{equation}
Every characteristic-two local ring with $\m^3=0$ is eligible; so is every cube-zero local ring with $2\m=0$ and even residue characteristic. Set $Q=|A|$, $S=\{a^2:a\in\m\}$ and $w=|S|/Q^2$.

\begin{lemma}[The square image]\label{lem:image}
$S$ is an ideal contained in $\m^2$, is annihilated by $\m$, and is a $k$-vector space. If $r=t+u$ with $u\in\m[x]$, then
\[
 r^2=t^2+u^2,\qquad u^2\in S[x^2].
\]
Conversely every polynomial in $S[x^2]$ is the square of some $u\in\m[x]$.
\end{lemma}
\begin{proof}
Since $2\m=0$, addition of squares on $\m$ is the square of a sum. The annihilation hypothesis makes the action of $A$ on $S$ factor through $k$. For $\lambda\in k$, take a square root in the finite field $k$ and lift it to $b\in A$. Then $\lambda a^2=(ba)^2$, showing that the image is stable under scalars. The cross terms in the polynomial square vanish because they have a factor $2$ multiplying an element of $\m$. Each coefficient in $S$ can be lifted independently to a square root in $\m$, giving the converse.
\end{proof}

Write $\bar f$ for reduction modulo $\m$ and $f_j$ for the coefficient of $x^j$. For $n\ge2$, define $D_n$ to consist of the $f\in B_n(A)$ with $0\le d=\deg\bar f\le n-2$ satisfying one of the following conditions (put $\bar f_{-1}=0$):
\begin{equation}\label{eq:D}
\begin{array}{ll}
 n-d\text{ even}: & (f_n,f_{n-1})=(s\bar f_d,s\bar f_{d-1})\text{ for some }s\in S;\\
 n-d\text{ odd}: & (f_n,f_{n-1})=(0,s\bar f_d)\text{ for some }s\in S.
\end{array}
\end{equation}
Products with residue-field elements in this formula use the canonical $k$-action on $S$, not a choice of coefficient field in $A$.

\begin{theorem}[Exact square-start characterization]\label{thm:exact}
For $f\in B_n(A)$ with $\bar f\ne0$, the following are equivalent:
\begin{enumerate}
\item $f\in D_n$;
\item there exists $r\in A[x]$ with $\deg\bar r\ge1$ and $r^2f\in B_n(A)$.
\end{enumerate}
If $\m^3=0$, these conditions are also equivalent to the existence of such an $r$ with both $rf$ and $r^2f$ in $B_n(A)$. The latter ratio can always be chosen with $\bar r=x$. Moreover
\begin{equation}\label{eq:count}
 \frac{|D_n|}{|B_n(A)|}=\frac{|S|}{Q^2}\left(1-q^{-(n-1)}\right).
\end{equation}
\end{theorem}
\begin{proof}
Suppose first that $z=r^2f$ belongs to the ball and $e=\deg\bar r\ge1$. Then $d+2e=\deg\bar z\le n$. Lift the coefficients of $\bar r$ to a polynomial $t$ of degree $e$ with unit leading coefficient $t_e$, and write $r=t+u$, $u\in\m[x]$. By Lemma~\ref{lem:image}, $z=t^2f+Uf$ with $U=u^2\in S[x^2]$ and $Uf=U\bar f$. Reduction modulo $S$ gives $z\bmod S=t^2(f\bmod S)$. Multiplication by a polynomial with unit leading coefficient adds its degree even over a ring with zero divisors. As $f\bmod S\ne0$, this shows $f_n,f_{n-1}\in S$.

If both top coefficients vanish, condition \eqref{eq:D} holds with $s=0$. Otherwise let $j=\deg f\in\{n-1,n\}$. Since $j+2e>n$, the leading terms of $t^2f$ and $U\bar f$ must cancel. In particular $U\ne0$ and
\[
 \deg U=j+2e-d.
\]
Here multiplication by the nonzero leading coefficient of $\bar f$ on the $k$-space $S$ is injective. The degree of $U$ is even, so $j-d$ is even. If $j=n-1$, then $n-d$ is odd, and the second line of \eqref{eq:D} holds because $\bar f_d\ne0$. If $j=n$, then $n-d$ is even. Let $v$ be the leading coefficient of $U$. The coefficients of $x^{n+2e}$ and $x^{n+2e-1}$ must both vanish; the latter degree exceeds $n$. The coefficient of $t^2$ one below its top is $2t_et_{e-1}$, whose product with $f_n\in S$ vanishes because $2\in\m$. There is no term of $U$ one below its top (all degrees are even). Thus in $S$,
\[
 \bar t_e^2 f_n+v\bar f_d=0,\qquad
 \bar t_e^2 f_{n-1}+v\bar f_{d-1}=0.
\]
Take $s=-v/\bar t_e^2$. This proves necessity, including the case of mixed characteristic.

Conversely assume \eqref{eq:D} and choose $c\in\m$ with $c^2=s$. Put
\[
 J=\begin{cases}(n+2-d)/2,&n-d\text{ even},\\
 (n+1-d)/2,&n-d\text{ odd},\end{cases}\qquad r_0=x+cx^J.
\]
Then $J$ is an integer, $J+d\le n$, and $r_0^2=x^2+sx^{2J}$. In the even case the two coefficients of $x^2f$ above degree $n$ cancel against the corresponding coefficients of $sx^{2J}\bar f$; in the odd case only the coefficient of degree $n+1$ needs cancellation. In both cases $2S=0$, so the plus sign cancels. There are no higher terms of $sx^{2J}\bar f$. Thus $r_0^2 f\in B_n(A)$.

If $\m^3=0$, every coefficient of $y_0=r_0f$ above degree $n$ belongs to $\m^2$: the contribution from $xf$ is $f_n\in S$, and $cx^J f$ has high-degree terms only from coefficients of $f$ in $\m$ since $J+d\le n$. The module $\m^2$ is annihilated by $\m$. Divide the part of $y_0$ above degree $n$ by $\bar f$ as a polynomial with coefficients in the $k$-vector space $\m^2$ (cancel leading coefficients successively). Its remainder has degree $<d$, or is zero if $d=0$. Let $v\in\m^2[x]$ be minus the quotient. For $r=r_0+v$ we have $rf=y_0+v\bar f$ of degree at most $n$, while $r^2=r_0^2$ because $v^2=0$ and $2v=0$. The reductions have degrees $d,d+1,d+2$ and are nonzero. This proves the exact triple-start claim.

Finally, a uniformly random polynomial in $B_n(A)$ has uniform independent residue coefficients. The probability of $0\le d\le n-2$ is $q^{-2}-q^{-(n+1)}$. Given any such residue polynomial, each top coefficient is independently uniform in $\m$; exactly $|S|$ pairs satisfy \eqref{eq:D}, since $\bar f_d$ is nonzero. Therefore the fraction is $(q^{-2}-q^{-(n+1)})|S|/(Q/q)^2$, as stated.
\end{proof}

\begin{remark}
When $\m^3=0$, this calculation is sharp in a stronger sense than an order-of-magnitude bound: among polynomials with nonzero reduction, $D_n$ is exactly the union of all possible first positions of triples with ratio of positive residue degree. It need not be an optimal hitting set, because one can instead remove some middle or last terms. Thus the theorem does not assert $\Delta_L(A)=1-w$.
\end{remark}

\section{A recoverable packing and local bounds}
\begin{lemma}[Deletion]\label{lem:deletion}
The set $B_n(A)\setminus(D_n\cup\{f:\bar f=0\})$ is admissible for every $L\ge2$ and has fraction at least $1-w-q^{-(n+1)}$.
\end{lemma}
\begin{proof}
A polynomial over a finite local ring is a unit exactly when its reduction is a nonzero constant. Indeed the ideal $\m[x]$ is nilpotent (the radical of a finite local ring is nilpotent), and a lift of a nonzero constant is a unit modulo that ideal. In an alleged pattern \eqref{eq:pattern}, all reductions are nonzero, so the nonunit ratio has positive residue degree. At least one of $a,b,a+b$ is even. The earlier term of the pair separated by that even difference is a square-start for a ratio of positive residue degree, hence lies in $D_n$ by Theorem~\ref{thm:exact}. This contradicts its retention. Counting follows from \eqref{eq:count}.
\end{proof}

\begin{lemma}[Packing]\label{lem:packing}
For $n\ge3$ there is a packing consisting of exactly $Q^{n+1}\gamma w$ pairwise vertex-disjoint triples in $B_n(A)$, where
\[
 \gamma=(1-q^{-1})^2.
\]
Each triple is $(f,rf,r^2f)$ with $\bar r=x$ and all three reductions nonzero.
\end{lemma}
\begin{proof}
Use all starts $f$ determined by arbitrary $f_0,\ldots,f_{n-2}$ with $f_0,f_{n-2}$ units, and $s\in S$, by setting
\[
 f_n=s f_{n-2},\qquad f_{n-1}=s f_{n-3}.
\]
Fix once and for all one square root $c(s)\in\m$ of each $s$. Define
\begin{equation}\label{eq:ratio}
 r=x+c(s)(x^2+1)+sx^3.
\end{equation}
Its square equals $x^2+s(x^4+1)$, since $2\m=0$, $s^2=0$ and $s\m=0$. In $z=r^2f$, the coefficients of degrees $n+2,n+1$ cancel and there are no higher terms. In $y=rf$, the only potential degree greater than $n$ is $n+1$, whose coefficient is $f_n+sf_{n-2}=0$: the products of $c$ or $s$ with the top two coefficients of $f$ are zero. Thus $y,z$ lie in the ball, and their residue degrees distinguish the three positions ($n-2,n-1,n$).

To show disjointness it suffices to recover the start from either $y$ or $z$. From $\bar y=x\bar f$ or $\bar z=x^2\bar f$ recover $\bar f_0\ne0$. The constant terms satisfy
\[
 y_0^2=s\bar f_0^2,\qquad z_0=s\bar f_0,
\]
so either one determines $s$ in the $k$-vector space $S$, hence the chosen $c(s)$ and the entire ratio $r$. Multiplication by $r$ is injective: if $rh=0$, reduce modulo $\m$ and obtain $x\bar h=0$, so $h\in\m[x]$; successively reduce modulo $\m^{j+1}$ to infer $h\in\m^{j+1}[x]$ whenever $h\in\m^j[x]$, since multiplication by $x$ is injective on $(\m^j/\m^{j+1})[x]$. Nilpotence concludes $h=0$. Multiplication by $r^2$ is also injective. Thus each position is injectively parametrized, and the different residue degrees separate the positions.

There are $(Q(1-1/q))^2 Q^{n-3}|S|$ choices, because $f_0$ and $f_{n-2}$ are distinct coefficients and $f_n$ determines $s$. This is $Q^{n+1}\gamma w$.
\end{proof}

\section{Products and passage to infinite sets}
For $k\ge0$, let $\alpha_k$ be the maximum fraction of the ternary cube $\{0,1,2\}^k$ containing no combinatorial line, with $\alpha_0=1$. A line fixes some coordinates and uses the same variable letter in every other coordinate, with at least one variable coordinate. The density Hales--Jewett theorem \cite{dhj} says $\alpha_k\to0$. Also $\alpha_k\le2/3$ for $k\ge1$ by slicing into lines parallel to a coordinate axis.

\begin{lemma}[Degree-block gluing]\label{lem:glue}
For every fixed finite product $R=\prod_{i=1}^h A_i$ of finite local rings, not necessarily eligible, and every fixed $L\ge2$,
\[
 \Delta_L(R)=\limsup_{n\to\infty}M_{L,n}(R).
\]
\end{lemma}
\begin{proof}
The inequality $\le$ follows by restriction. For the reverse, choose cutoffs $n_j$ approaching the upper limit of the finite optima and growing so fast that, with $n_0=0$,
\[
 n_j-j>(1+L)n_{j-1},\qquad n_j-(1+L)j>n_{j-1}.
\]
At cutoff $n_j$, retain from an optimizer only polynomials for which every residue degree is in $[n_j-j,n_j]$. The discarded fraction is at most $\sum_i q_i^{-j-1}$ (including the zero residues). Take the union of all retained blocks. In a putative pattern, the residue of the ratio is nonzero in every coordinate since all three polynomials have nonzero reductions. In each coordinate the residue degrees form $d,d+ae,d+(a+b)e$, with $e\ge0$; hence block indices are nondecreasing. If all terms are in one block, admissibility is already known. Otherwise consider the latest block $j$. If only the last term is in this block, in any coordinate the preceding two degrees are at most $n_{j-1}$, so $e\le n_{j-1}$ and the last degree is at most $(1+L)n_{j-1}$, contradiction. If the last two terms are in block $j$, their degree difference is at most $j$; hence $e\le j$ and the first degree is at least $n_j-(1+L)j$, contradiction. This proves global admissibility. The fractions along $n_j$ approach the finite upper limit from below, proving equality.
\end{proof}

\begin{theorem}[Product-density classification]\label{thm:main}
Let $(A_i,\m_i)$ be eligible finite local rings. Write $q_i=|A_i/\m_i|$, $Q_i=|A_i|$, $S_i=\{a^2:a\in\m_i\}$, $w_i=|S_i|/Q_i^2$, and $\gamma_i=(1-q_i^{-1})^2$. For any finite $h$ and $L\ge2$,
\begin{equation}\label{eq:product}
 \prod_{i=1}^h(1-w_i)\le \Delta_L\left(\prod_{i=1}^h A_i\right)
 \le \E\alpha_K,\quad K=\sum_{i=1}^h X_i,
\end{equation}
where $X_i$ are independent Bernoulli variables of success probabilities $3\gamma_i w_i$. In particular,
\begin{equation}\label{eq:local}
 \gamma_i w_i\le 1-\Delta_L(A_i)\le w_i.
\end{equation}
For an infinite sequence of factors the limit of $\Delta_L(\prod_{i=1}^h A_i)$ as $h\to\infty$ exists and is strictly positive if and only if $\sum_{i\ge1}w_i<\infty$.
\end{theorem}
\begin{proof}
At finite cutoff, take the Cartesian product of the sets in Lemma~\ref{lem:deletion}. If a pattern with nonunit ratio lies in the product then at least one coordinate has nonunit ratio, since units in a finite product are componentwise units. All residue terms are nonzero. In this coordinate the residue degree of the ratio is positive and the terms are distinct, contradicting deletion. The product fraction is at least $\prod_i(1-w_i-q_i^{-(n+1)})$ for sufficiently large $n$. Apply Lemma~\ref{lem:glue} and let $n\to\infty$.

For the upper bound, partition each $B_n(A_i)$ into the triples in Lemma~\ref{lem:packing} and singleton cells. A uniform point lands in a triple cell with probability $3\gamma_iw_i$ independently in each coordinate. Condition on the chosen cells. The product cell is a ternary cube in $K$ coordinates (with singleton coordinates suppressed). A combinatorial line consists of a three-term geometric progression in the product: use the ratio of the chosen triple on every moving coordinate, and $1$ elsewhere. At least one coordinate moves, so the ratio is a nonunit and the terms are distinct. An admissible subset occupies at most $\alpha_K$ of the cell. Average over cells, giving the upper bound at every $n\ge3$ and thus for $\Delta_L$. For $h=1$, use $\alpha_1=2/3$ to obtain \eqref{eq:local}.

By slicing finite optimizers along the last coordinate, $M_{L,n}(\prod_{i=1}^{h+1} A_i)\le M_{L,n}(\prod_{i=1}^h A_i)$: a forbidden pattern in a slice extends by the ratio $1$ in its last coordinate. Lemma~\ref{lem:glue} implies the same monotonicity for $\Delta_L$. Thus a limit exists. Since $S_i\subseteq\m_i$, $w_i\le (q_iQ_i)^{-1}\le1/4$. If $\sum_i w_i<\infty$, the infinite product of $1-w_i$ is positive, by $-\log(1-w_i)\le(4/3)w_i$. If $\sum_i w_i=\infty$, use $\gamma_i\ge1/4$ to obtain $\mu_h=\E K=3\sum_{i\le h}\gamma_iw_i\to\infty$. Independence gives $\operatorname{Var}K\le\mu_h$, and Chebyshev's inequality gives
\[
 \E\alpha_K\le 4/\mu_h+\sup_{k\ge\mu_h/2}\alpha_k\longrightarrow0.
\]
This proves both directions.
\end{proof}

\begin{corollary}[Uniform endpoints without nesting]\label{cor:endpoints}
For any sequence of finite products of eligible rings, put $W=\sum_i w_i$ for each product. For fixed $L$, $\Delta_L\to0$ if and only if $W\to\infty$, and $\Delta_L\to1$ if and only if $W\to0$. In the small-loss regime, with $q_* =\min_i q_i$,
\[
 1-\Delta_L=W+O(W^2+W/q_*),
\]
with an absolute implied constant.
\end{corollary}
\begin{proof}
The proof of Theorem~\ref{thm:main} is uniform in the factors. In addition to the estimate using $\mu=\E K\ge3W/4$, its lower bound gives $\Delta_L\ge e^{-4W/3}$, and
\[
 \frac{1-\prod_i(1-3\gamma_iw_i)}{3}\le1-\Delta_L\le W.
\]
The left side is at least $(1-e^{-3W/4})/3$. These inequalities give both equivalences, by considering subsequences if a stated condition fails. For $W$ small, the left side equals $\sum_i\gamma_iw_i+O(W^2)$ by the first two inclusion--exclusion inequalities; $0\le1-\gamma_i\le2/q_i$ gives the asserted uniform expansion.
\end{proof}

\section{Same size and radical layers, opposite limits}
Suppose $A$ has characteristic two and $\m^3=0$. Then the squaring map
\[
 \m/\m^2\longrightarrow\m^2,\qquad v\longmapsto v^2
\]
is well-defined, additive and Frobenius-semilinear over $k$. Its image is $S$. If $r$ is the dimension of that image over $k$, the invariant is
\begin{equation}\label{eq:rank}
 w=q^r/Q^2.
\end{equation}
Indeed $\m^2$ is annihilated by $\m$ and its square is zero, so the square map factors through the indicated quotient and scalar multiplication is compatible with the Frobenius automorphism. In particular the dimensions of the two radical layers alone do not determine $w$.

\begin{theorem}[Opposite limits with indistinguishable layer sizes]\label{thm:opposite}
For every fixed power $q$ of two, there exist sequences $(A_i)$ and $(C_i)$ of characteristic-two local $\F_q$-algebras, with $\m^3=0\ne\m^2$ for every factor, such that corresponding factors have the same cardinality, residue field, and dimensions of all radical layers, but for every fixed $L\ge2$,
\[
 \lim_{h\to\infty}\Delta_L(A_1\times\cdots\times A_h)>0,
 \qquad
 \lim_{h\to\infty}\Delta_L(C_1\times\cdots\times C_h)=0.
\]
\end{theorem}
\begin{proof}
For each integer $r\ge1$ take a vector space basis $1,x_1,\ldots,x_r,y_1,\ldots,y_r,z_1,\ldots,z_r$. In $A(r)$ define $x_i y_i=y_i x_i=z_i$ and let all other products of nonidentity basis elements vanish. In $C(r)$ define $x_i^2=z_i$ and let all other such products vanish. Both are commutative associative algebras: any triple product of radical elements vanishes, so associativity reduces to the identity cases and bilinearity. The radical is the span of all basis elements except $1$; it is nilpotent, and an element is a unit exactly when its constant coefficient is nonzero. In both rings the radical square is the span of the $z_i$ and the radical cube is zero. The dimensions of $\m/\m^2$ and $\m^2$ are $2r$ and $r$, respectively, and $Q=q^{1+3r}$ for both.

In characteristic two, every radical element in $A(r)$ squares to zero: the cross terms vanish even though $x_i y_i$ need not vanish. In $C(r)$, the image of squaring is the full span of the $z_i$, by perfection of $\F_q$. Thus
\[
 w(A(r))=q^{-2-6r},\qquad w(C(r))=q^{-2-5r}.
\]
In both sequences use $q^{5r}$ copies of the ring of rank $r$, in successive blocks for $r=1,2,\ldots$. The parameter sums are respectively
\[
 \sum_{r\ge1}q^{5r}q^{-2-6r}=\frac{q^{-2}}{q-1},
 \qquad
 \sum_{r\ge1}q^{5r}q^{-2-5r}=\sum_{r\ge1}q^{-2}=\infty.
\]
Apply Theorem~\ref{thm:main}. In fact the positive limit for the first sequence is at least $\exp[-4/(3q^2(q-1))]$.
\end{proof}

\section{Scope, numerical checks, and next questions}
\paragraph{What is proved.} Theorem~\ref{thm:main} supplies a necessary-and-sufficient criterion for positive limiting density throughout the eligible class, not just an obstruction or a reduction. Theorem~\ref{thm:exact} computes the exact local obstruction and characterizes possible triple starts for cube-zero radicals. Theorem~\ref{thm:opposite} shows that the multiplication law, more precisely its squaring image, enters the criterion even with every radical-layer size fixed. This is a partial answer to Asada et al.'s published Question 6.5 and a resolution of the corresponding characteristic-two cube-zero subproblem posed in the repository. It is not a proof of the full question for arbitrary finite rings, nor a claimed resolution of a major outstanding number-theoretic conjecture such as Artin's primitive root conjecture or Chowla's central nonvanishing conjecture.

\paragraph{Relation to earlier bounds.} In the square-zero case $S=0$ as a vector space, $|S|=1$, recovering the earlier $Q^{-2}$ parameter; the earlier chapter gives a sharper single-factor constant in that case. Our new packing has a uniform constant $\gamma\ge1/4$ that works also for nonzero square image. More generally the hypotheses allow rings of larger nilpotence length if all radical squares are annihilated by the radical (for example, the commutative characteristic-two algebra with generators $u_i$ subject to $u_i^2=0$); the square-pair and product theorems still hold there. The exact triple-start converse is only proved here when $\m^3=0$.

\paragraph{Exact regression tests.} The included pure-Python script checks all degree-$\le3$ polynomials in six small rings. It enumerates every relevant squared ratio (with a degree cap justified by the leading term comparison), tests the exact criterion for every nonzero-reduction polynomial, constructs the radical-square correction for every obstructed start, and tests vertex-disjointness of the displayed packing. Results are in Table~\ref{tab:check}. No numerical computation is used in the theorems.
\begin{table}[ht]
\centering\small
\begin{tabular}{lrrrr}
\toprule
Ring & $Q$ & $|S|$ & $|D_3|$ & packed triples\\
\midrule
$\F_2$ & 2 & 1 & 3 & 1\\
$\mathbb Z/4\mathbb Z$ & 4 & 1 & 12 & 4\\
$\F_2[t]/(t^3)$ & 8 & 2 & 96 & 32\\
$A(1)$ over $\F_2$ & 16 & 1 & 192 & 64\\
$C(1)$ over $\F_2$ & 16 & 2 & 384 & 128\\
$(\mathbb Z/4\mathbb Z)[t]/(2t,t^3)$ & 16 & 2 & 384 & 128\\
\bottomrule
\end{tabular}
\caption{Exhaustive checks at cutoff $n=3$; each deletion count equals the number of square-starts. The script also verifies the triple-start construction.}\label{tab:check}
\end{table}

\paragraph{Open directions.} The sharp value of $\Delta_L(A)$ for a single eligible ring is not determined. A sharper packing may change the factor $(1-q^{-1})^2$ but cannot change the product summability threshold. When the residue characteristic is two but $2\m\ne0$, the square contains a cross term $2tu$; when $\m S\ne0$, the replacement $Uf=U\bar f$ is invalid. In odd residue characteristic the even-gap elimination no longer gives our square formula. Each of these requires a new argument. No Lean kernel check of either the upstream development or this note has been performed; the mathematical proof here is self-contained except for the classical density Hales--Jewett theorem and elementary facts about finite local rings. Novelty is relative to the inspected merged article and a targeted literature search, not a guarantee of global priority.

\begin{thebibliography}{9}
\bibitem{asada} M. Asada, E. Fourakis, S. Manski, N. McNew, S. J. Miller and G. Moreland, Subsets of $\mathbb F_q[x]$ free of 3-term geometric progressions, \emph{Finite Fields Appl.} \textbf{44} (2017), 135--147. \href{https://doi.org/10.1016/j.ffa.2016.10.002}{doi:10.1016/j.ffa.2016.10.002}; \href{https://arxiv.org/abs/1512.01932}{arXiv:1512.01932} (Question 6.5).
\bibitem{dhj} D. H. J. Polymath, A new proof of the density Hales--Jewett theorem, \emph{Ann. of Math.} \textbf{175} (2012), 1283--1327. \href{https://doi.org/10.4007/annals.2012.175.3.6}{doi:10.4007/annals.2012.175.3.6}.
\bibitem{proveit} ProveIt, \emph{Beyond the Riemann hypothesis for Dirichlet $L$-functions: consequences, transfers and limits}, repository article, 27 November 2026 snapshot. \href{https://github.com/VladimirReshetnikov/ProveIt/tree/main/Analysis/GeneralizedRiemannHypothesis}{Project tree}; \href{https://github.com/VladimirReshetnikov/ProveIt/blob/main/Analysis/GeneralizedRiemannHypothesis/Article/sections/gpp.tex}{density section} (blob \texttt{29c75c077d43ddcea05f9d68fd47e569c33ee69e}), labels \texttt{gpp:thm:local}, \texttt{gpp:q:cube}.
\bibitem{grh} L. Sernau et al., \emph{grh-lean}, pinned upstream revision \texttt{72fead102f663d5c7a407964a13dc0cf9d977207}, \href{https://github.com/lsernau/grh-lean}{repository}. Source declarations inspected in ProveIt's vendored copy; see accompanying source audit.
\end{thebibliography}
\end{document}
